\documentclass[10pt,a4paper]{amsart}
\usepackage[margin=19mm]{geometry}
\usepackage{amsmath,amssymb,mathtools}
\usepackage[scr=rsfs,cal=euler]{mathalfa}
\usepackage{xfrac}
\usepackage[unicode,hidelinks,bookmarks=false]{hyperref}
\newcommand{\ie}{i.e.\ }
\newcommand{\cf}{cf.\ }
\newcommand{\resp}{respectively\ }
\newcommand{\Subsection}{\subsection}
\providecommand{\nobreakdash}{\nobreak\hbox{-}}
\setlength{\emergencystretch}{2em}
\multlinegap0pt
\usepackage[shortlabels]{enumitem}
\usepackage{xcolor}
\usepackage{tcolorbox}
\usepackage{algorithm}\usepackage{algpseudocode}
\usepackage{amssymb}
\usepackage{mathtools}
\newtheorem{lemma}{Lemma}
\title{Black-Box Detection of LLM-Generated Text Using Generalized Jensen-Shannon Divergence}
\author{Shuangyi Chen, Ashish Khisti}
\date{Source: arXiv:2510.07500v3 (CC BY 4.0)}
\begin{document}
\maketitle

\noindent\emph{Source and licence.} Black-Box Detection of LLM-Generated Text Using Generalized Jensen-Shannon Divergence. Version arXiv:2510.07500v3 (CC BY 4.0), distributed by arXiv under the Creative Commons Attribution 4.0 International licence, http://creativecommons.org/licenses/by/4.0/, as recorded in the arXiv licence metadata for this exact version and shown on the abstract page https://arxiv.org/abs/2510.07500v3. Full licence notice: \texttt{licenses/arxiv-2510.07500-CC-BY-4.0.txt}.\par\smallskip

\noindent\textit{Editorial scope.} Counted unit: the article's lemma lem:uniform-bound (source0.tex lines 1466--1474). The article's complete original proof of it is delivered with it (source0.tex lines 1476--1514, one \texttt{proof} environment of 39 lines). No mathematics is written, repaired or re-derived: the statement and the proof are the article's own lines, in the article's own notation, and no retained line is substituted or re-flowed.  No further source-local context is needed: every cross-reference of every retained line resolves inside the two delivered spans.  Nothing was omitted from any retained span, so the package carries no omission record.  No retained line contains a citation, so the wrapper carries no bibliography and no bibliography entry.  No retained line contains a figure environment, an image-inclusion command or drawing code, so no image is delivered and the package declares no asset dependency.  Editorial formatting only, in addition: an AMS \texttt{amsart} preamble that restates the article's own declarations and macros used by the retained lines, namely the environments \texttt{lemma} on separate counters, together with the standard packages those lines need.  The retained environments print in this document as Lemma 1 in order of appearance, whereas the article prints the same items with the numbering of its own full document.  The complete native source of the article as distributed in the arXiv e-print for this exact version is bundled unchanged as \texttt{sources/186678/source0.tex}.
\begin{lemma}
\label{lem:uniform-bound}
For all $\alpha > 0$ and $p \in (0,1]$, it holds that
\begin{equation}
  p \, \max\{1, \log(1/p)\} \, e^{-\alpha p}
  \;\;\le\;\; \frac{2 + \log(1+\alpha)}{e\,\alpha}.
  \label{eq:uniform-ineq}
\end{equation}
\end{lemma}

\begin{proof}
Let $y = \alpha p \in (0,\alpha]$ and $A = \log \alpha$. Then we can rewrite
\[
p \, \max\{1,\log(1/p)\} e^{-\alpha p}
= \frac{1}{\alpha} \, y e^{-y} \, \max\{1, A - \log y\}.
\]
Next observe the inequality
\[
\max\{1, A-\log y\} \;\le\; 1 + A_{+} + (-\log y)_{+},
\]
where $x_{+}=\max\{0,x\}$ and $A_{+}=\max\{0,A\}$.

Therefore,
\[
y e^{-y} \max\{1,A-\log y\}
\;\le\; (1+A_{+}) \cdot y e^{-y} + y e^{-y}(-\log y)_{+}.
\]

Now use the following standard bounds:
\[
\sup_{y>0} y e^{-y} = \frac{1}{e}, 
\qquad
\sup_{0<y\le 1} y(-\log y) = \frac{1}{e}.
\]
Hence
\[
\sup_{y>0} y e^{-y} \max\{1, A-\log y\}
\;\le\; \frac{1+A_{+}}{e} + \frac{1}{e}
= \frac{2 + \log\alpha_{+}}{e},
\]
where $\log\alpha_{+} = \max\{0,\log \alpha\} \le \log(1+\alpha)$.

Substituting back into the expression, we obtain
\[
p \, \max\{1, \log(1/p)\} e^{-\alpha p}
\le \frac{1}{\alpha}\cdot \frac{2+\log(1+\alpha)}{e}.
\]
This proves Eq.~\ref{eq:uniform-ineq}.
\end{proof}

\end{document}
